\documentclass[11pt,a4paper]{article}

\usepackage[T1]{fontenc}
\usepackage[utf8]{inputenc}
\usepackage{lmodern}
\usepackage{microtype}
\usepackage[margin=2.3cm,top=2.6cm,bottom=2.4cm]{geometry}
\usepackage{booktabs}
\usepackage{tabularx}
\usepackage{array}
\usepackage[table]{xcolor}
\usepackage{titlesec}
\usepackage{enumitem}
\usepackage{fancyhdr}
\usepackage{caption}
\usepackage{amssymb}
\usepackage{graphicx}
\usepackage[colorlinks=true,linkcolor=black,urlcolor=black,citecolor=black]{hyperref}

\definecolor{ink}{HTML}{000000}
\color{ink}

\titleformat{\section}
  {\normalfont\Large\bfseries}{\thesection}{0.7em}{}[\vspace{-0.55em}{\color{black}\rule{\linewidth}{0.7pt}}]
\titlespacing*{\section}{0pt}{1.7em}{0.7em}
\titleformat{\subsection}{\normalfont\large\bfseries}{\thesubsection}{0.6em}{}
\titlespacing*{\subsection}{0pt}{1.3em}{0.5em}
\titleformat{\subsubsection}{\normalfont\normalsize\bfseries}{\thesubsubsection}{0.6em}{}

\pagestyle{fancy}
\setlength{\headheight}{22pt}
\addtolength{\topmargin}{-10pt}
\fancyhf{}
\renewcommand{\headrulewidth}{0.3pt}
\fancyhead[L]{\small Submission 3 — Custom CNN from Scratch}
\fancyhead[R]{\small Cricket Shot Attention CNN}
\fancyfoot[C]{\small\thepage}

\newenvironment{keybox}
  {\par\medskip\noindent\begin{tabular}{@{}p{0.3em}@{\hspace{0.8em}}p{\dimexpr\linewidth-2.2em}@{}}
   {\color{black}\vrule width 2.5pt} &}
  {\end{tabular}\par\medskip}

\setlist[itemize]{leftmargin=1.3em,itemsep=0.25em,topsep=0.4em}
\setlist[enumerate]{leftmargin=1.5em,itemsep=0.3em,topsep=0.4em}
\renewcommand{\arraystretch}{1.25}
\captionsetup{font=small,labelfont=bf,skip=6pt}
\setlength{\parskip}{0.55em}
\setlength{\parindent}{0pt}
\graphicspath{{figures/}}

\begin{document}

\begin{center}
{\small CNN PROJECT \textbullet\ SUBMISSION 3 OF 4}\\[0.9em]
{\LARGE\bfseries Custom CNN Trained from Scratch}\\[1.0em]
{\large Four convolutional blocks, 1.2 M parameters, on a leakage-controlled split}\\[1.2em]
{\small Repository: \href{https://github.com/Rayhan-who/cricket-shot-attention-cnn}{github.com/Rayhan-who/cricket-shot-attention-cnn}}
\end{center}

\vspace{0.4em}{\hrule height 0.8pt}\vspace{1.0em}

\section*{Summary}

A four-block convolutional network with 1,206,756 trainable parameters was trained from
random initialisation for 50 epochs. It reached \textbf{0.8514} test accuracy and
\textbf{0.8515} macro F1 on a 666-image held-out test set, in 29.9 minutes on a Kaggle
T4 GPU.

Three results matter more than the headline figure.

\begin{enumerate}
\item The network \textbf{did not overfit and did not converge}. The
  train--validation accuracy gap at the final epoch was $+0.0042$, and validation loss
  reached its minimum at epoch 50 and was still falling. Early stopping never triggered.
\item The per-class difficulty ordering \textbf{matches Submission 2's mean ordering},
  across four architectures and two different data splits.
\item The hardest confused pair is \texttt{drive} $\leftrightarrow$
  \texttt{legglance-flick}, which \textbf{confirms the visual prediction made in
  Submission 1 and overturns the measurement made in Submission 2}.
\end{enumerate}

\begin{keybox}
\textbf{This run was made possible by a leakage audit conducted between Submissions 2
and 3}, which established that 45.9\% of Submission 2's test set contained
near-duplicates of training images. All results here use a corrected split. The audit is
summarised in Section~\ref{sec:audit}.
\end{keybox}

\section{Architecture}

\begin{tabularx}{\linewidth}{@{}>{\ttfamily}p{3.6cm}p{4.4cm}X@{}}
\toprule
\textnormal{\bfseries Stage} & \textbf{Composition} & \textbf{Output} \\
\midrule
input       & --- & $224 \times 224 \times 3$ \\
augmentation & rotation, zoom, brightness, contrast & $224 \times 224 \times 3$ \\
rescale     & $\times\,1/255$ & $224 \times 224 \times 3$ \\
block 1     & (Conv$3{\times}3$ $\to$ BN $\to$ ReLU) $\times 2$, MaxPool, SpatialDropout & $112 \times 112 \times 32$ \\
block 2     & same, 64 filters  & $56 \times 56 \times 64$ \\
block 3     & same, 128 filters & $28 \times 28 \times 128$ \\
block 4     & same, 256 filters & $14 \times 14 \times 256$ \\
gap         & global average pooling & 256 \\
fc1         & Dense $\to$ BN $\to$ ReLU $\to$ Dropout & 128 \\
predictions & Dense, softmax & 4 \\
\bottomrule
\end{tabularx}

Eight convolutional layers in total. \textbf{1,206,756} trainable parameters --- less
than a third of EfficientNetB0's 3.97 M and about one twentieth of ResNet50's 23.5 M.

\subsection{Design decisions}

\begin{tabularx}{\linewidth}{@{}p{4.6cm}X@{}}
\toprule
\textbf{Choice} & \textbf{Reason} \\
\midrule
BatchNorm before every ReLU & Training from random initialisation on 3,361 images is unstable without it \\
\textbf{BatchNorm momentum 0.90} & Keras defaults to 0.99, which needs $\sim$200 updates for the running statistics to converge. On a dataset this size validation is then scored with un-converged statistics and appears broken for the first 20 epochs --- see Section~\ref{sec:bn} \\
Two convolutions per block & A $3{\times}3$ pair has the receptive field of a $5{\times}5$ with fewer parameters and an extra nonlinearity \\
Global average pooling & Flattening $14{\times}14{\times}256$ into Dense(128) would add 6.4 M parameters, more than the rest of the network combined \\
Spatial dropout 0.10 per block & Drops whole feature maps rather than individual pixels, the correct form for convolutional layers \\
Dropout 0.40 in the head & The densest part of the network and the easiest place to overfit \\
Rescaling $1/255$ & No ImageNet statistics apply; there is no pretrained backbone to match \\
\bottomrule
\end{tabularx}

The convolutional block is a single function taking an optional \texttt{attention}
argument, unused here. Submission 4 passes a CBAM module into that argument and changes
nothing else, which is what makes the comparison attributable.

\subsection{A note on BatchNorm momentum}
\label{sec:bn}

An initial three-epoch run produced a misleading picture: training accuracy climbed
normally from 0.348 to 0.433 while validation loss rose monotonically from 1.641 to
1.849, and the notebook's own diagnostic declared the model to be overfitting.

It was not. BatchNorm uses batch statistics during training but \emph{running averages}
at validation time. At the Keras default momentum of 0.99, those averages require on the
order of 200 updates to become meaningful. With 106 steps per epoch, validation for the
first several epochs was being computed with effectively broken normalisation.

Lowering momentum to 0.90 resolved it completely. In the run reported here, validation
tracks training from the first epoch onward (Figure~\ref{fig:curves}). This is recorded
because it is a failure mode that looks exactly like overfitting and leads directly to
the wrong corrective action --- adding regularisation to a model that needs none.

\section{Training}

\begin{tabularx}{\linewidth}{@{}>{\bfseries}p{5.0cm}X@{}}
\toprule
Setting & Value \\
\midrule
Split & clean (embedding clusters): 3,361 train / 697 validation / 666 test \\
Seed & 42 \\
Batch size & 32 (106 steps per epoch) \\
Optimizer & Adam, learning rate $10^{-3}$ \\
Epoch budget & 50 \\
Early stopping & validation loss, patience 12, not before epoch 5 \\
Learning-rate schedule & ReduceLROnPlateau, factor 0.3 \\
Augmentation & rotation $\pm10^\circ$, zoom 0.8--1.0, brightness $\pm20\%$, contrast $\pm20\%$ \\
Horizontal flip & \textbf{not applied} --- mirroring inverts the label on leg-side strokes \\
Class weights & not applied --- imbalance is 1.12\,:\,1 \\
\bottomrule
\end{tabularx}

The rubric requires a minimum of five epochs; \texttt{start\_from\_epoch=5} enforces this
regardless of how the validation curve behaves.

\section{Results}

\begin{center}
\begin{tabular}{@{}lr@{}}
\toprule
\textbf{Metric} & \textbf{Value} \\
\midrule
Test accuracy & 0.8514 \\
Precision (macro) & 0.8506 \\
Recall (macro) & 0.8536 \\
\textbf{F1 (macro)} & \textbf{0.8515} \\
F1 (weighted) & 0.8510 \\
\midrule
Epochs run & 50 / 50 \\
Training time & 29.87 min (35.8 s per epoch) \\
Testing time & 1.95 s (2.93 ms per image) \\
Trainable parameters & 1,206,756 \\
Train--validation gap & $+0.0042$ \\
\bottomrule
\end{tabular}
\end{center}

\subsection{Training curves}

\begin{figure}[htbp]
\centering
\includegraphics[width=\linewidth]{06_s3_curves.png}
\caption{Accuracy and loss, training versus validation, over 50 epochs. The two curves
track each other closely throughout, and both are still improving at the budget limit.
The transient spike in validation loss around epoch 25 coincides with a learning-rate
reduction and resolves within three epochs.}
\label{fig:curves}
\end{figure}

\begin{keybox}
\textbf{The model is budget-limited, not capacity-limited.} Validation loss reached its
minimum at epoch 50 --- the last epoch --- and early stopping never fired. A longer
schedule would score higher. Submission 4 inherits the identical 50-epoch budget, so the
comparison remains controlled, but both models should be described as compared
\emph{under a fixed budget} rather than \emph{at convergence}.
\end{keybox}

\subsection{Confusion matrix and per-class metrics}

\begin{figure}[htbp]
\centering
\includegraphics[width=\linewidth]{07_s3_confusion.png}
\caption{Test-set confusion matrix (counts, row-normalised shading) and per-class
precision, recall and F1.}
\end{figure}

\begin{center}
\begin{tabular}{@{}lrrrrl@{}}
\toprule
\textbf{Class} & \textbf{Precision} & \textbf{Recall} & \textbf{F1} & \textbf{Support} & \\
\midrule
\texttt{sweep}           & 0.9222 & 0.9432 & 0.9326 & 176 & easiest \\
\texttt{pullshot}        & 0.8345 & 0.8897 & 0.8612 & 136 & \\
\texttt{legglance-flick} & 0.8788 & 0.8056 & 0.8406 & 180 & \\
\texttt{drive}           & 0.7670 & 0.7759 & 0.7714 & 174 & hardest \\
\bottomrule
\end{tabular}
\end{center}

\subsubsection*{The ordering is a property of the data, not of the architecture}

\begin{center}
\begin{tabular}{@{}lcc@{}}
\toprule
\textbf{Class} & \textbf{Rank, Submission 3} & \textbf{Rank, Submission 2} \\
\midrule
\texttt{sweep}           & 1st & 1st \\
\texttt{pullshot}        & 2nd & 2nd \\
\texttt{legglance-flick} & 3rd & 3rd \\
\texttt{drive}           & 4th & 4th \\
\bottomrule
\end{tabular}
\end{center}

This ordering matches Submission 2's mean ordering across six runs (three pretrained
architectures, two transfer-learning modes). In Submission 2 the extremes were stable
--- \texttt{sweep} easiest and \texttt{drive} hardest in five of six runs --- while the two
middle classes swapped between runs. That the same ordering reappears in a from-scratch
network on a different data split points to a property of the classes themselves rather
than of any particular model.

The explanation is geometric. \texttt{sweep} is the only stroke played from a lowered or
kneeling position, so its silhouette differs from the other three in a way that survives
downsampling to $224 \times 224$. \texttt{drive} is an upright, predominantly side-on
stroke that shares its body geometry with both \texttt{legglance-flick} and
\texttt{pullshot}, leaving bat angle at contact as almost the only discriminating cue.

\subsection{The hardest pair, and a correction to Submission 2}

\begin{center}
\begin{tabular}{@{}lll@{}}
\toprule
& \textbf{Hardest pair} & \textbf{Basis} \\
\midrule
Submission 1 & \texttt{drive} $\leftrightarrow$ \texttt{legglance-flick} & human inspection of the sample grid \\
Submission 2 & \texttt{drive} $\leftrightarrow$ \texttt{pullshot} & confusion matrix, \textbf{contaminated} split \\
\textbf{Submission 3} & \textbf{\texttt{drive} $\leftrightarrow$ \texttt{legglance-flick}} & confusion matrix, \textbf{clean} split \\
\bottomrule
\end{tabular}
\end{center}

Counting both directions on the clean split:

\begin{center}
\begin{tabular}{@{}lrr@{}}
\toprule
\textbf{Pair} & \textbf{Errors} & \textbf{Share of 99 total} \\
\midrule
\texttt{drive} $\leftrightarrow$ \texttt{legglance-flick} & \textbf{40} & 40\% \\
\texttt{drive} $\leftrightarrow$ \texttt{pullshot} & 27 & 27\% \\
\bottomrule
\end{tabular}
\end{center}

Submission 2 identified the wrong pair because its best model made only \emph{two} errors
in 721 predictions --- a confusion matrix built on two data points carries no
information. That near-perfect score was itself a product of the contamination.

\begin{keybox}
\textbf{The visual judgement made in Submission 1 was correct, and the contaminated
measurement in Submission 2 was not.} This is a concrete demonstration of how leakage
distorts not only headline accuracy but the secondary analysis built on top of it. The
\texttt{drive}/\texttt{legglance-flick} confusion is the specific weakness the attention
module in Submission 4 must address.
\end{keybox}

\section{Comparison with Submission 2}

\begin{figure}[htbp]
\centering
\includegraphics[width=\linewidth]{08_s3_vs_pretrained.png}
\caption{The from-scratch network against all six pretrained configurations.
\textbf{The bars are not measured on the same data} --- see the warning below.}
\end{figure}

\begin{keybox}
\textbf{This comparison is not like-for-like and must not be reported as though it
were.} Submission 2's figures were measured on the hash-grouped split, where 45.9\% of
the test set had a near-duplicate in training. This submission's figure is measured on
the clean split. The apparent 0.1458 macro-F1 gap therefore conflates two separate
effects: the genuine cost of training from scratch, and the removal of leakage. Nothing
in the present data separates them.
\end{keybox}

\subsection*{The experiment that would separate them}

Re-running a single pretrained model --- EfficientNetB0 fine-tuned --- on the clean
split would take approximately eleven minutes and would yield the two numbers needed:

\begin{center}
\begin{tabular}{@{}lcc@{}}
\toprule
& \textbf{Contaminated split} & \textbf{Clean split} \\
\midrule
EfficientNetB0 fine-tuned & 0.9973 & \emph{not yet measured} \\
Custom CNN from scratch & \emph{not measured} & 0.8515 \\
\bottomrule
\end{tabular}
\end{center}

Filling the top-right cell gives both the true cost of leakage for a pretrained model and
a valid like-for-like transfer-learning comparison. It is the single most informative
experiment still outstanding in this project.

\section{Headroom for Submission 4}

\begin{tabularx}{\linewidth}{@{}>{\bfseries}p{6.4cm}X@{}}
\toprule
Quantity & Value \\
\midrule
Baseline macro F1 & 0.8515 \\
Test set size & 666 \\
Standard error at this accuracy & 1.38 pp \\
Run-to-run variance (seed, shuffling, augmentation) & $\approx$ 0.70 pp \\
\textbf{Minimum detectable gain} & \textbf{$\approx$ 3.1 pp} \\
\textbf{Headroom available} & \textbf{14.8 pp} \\
\bottomrule
\end{tabularx}

The comparison is viable but not generously so. CBAM typically delivers 1--3 pp on tasks
of this kind, which sits at the detection threshold. If the measured gain falls below
3 pp, the defensible responses are to run both models with three seeds --- lowering the
threshold to about 1.8 pp --- or to report the result as inconclusive. Reporting a 1 pp
gain from a single run as evidence that attention helps would not be supportable.

\begin{keybox}
\textbf{For contrast: had Submission 2's contaminated 0.9973 been used as the baseline,
the gain required to clear noise would have exceeded the total headroom available, and
Submission 4 could not have produced a result at all.} The leakage audit is what made
this submission measurable.
\end{keybox}

\section{The leakage audit}
\label{sec:audit}

Conducted between Submissions 2 and 3, and the reason this submission uses a different
split from its predecessor.

\subsection{Motivation}

Submission 2's best model reached 0.9972 accuracy --- 719 of 721 test images correct ---
on a task characterised in Submission 1 as fine-grained. Submission 1 had controlled for
duplication by grouping images with identical perceptual hashes, removing 608 redundant
files (12.9\%). That control proved incomplete.

\subsection{Method}

Every image was passed through EfficientNetB0 with ImageNet weights as a fixed feature
extractor, producing a 1,280-dimensional vector. Vectors were L2-normalised and clustered
by connected components over cosine similarity. The threshold was swept from 0.99 to
0.90.

The rationale: a perceptual hash is computed from greyscale brightness gradients, so any
transformation that alters the channel mix alters the hash. A convolutional feature
extractor trained with colour augmentation is largely invariant to exactly those
transformations.

\subsection{Findings}

\begin{center}
\begin{tabular}{@{}rrrrr@{}}
\toprule
\textbf{Threshold} & \textbf{Clusters} & \textbf{Redundant} & \textbf{\%} & \textbf{Largest cluster} \\
\midrule
0.99 & 4,057 & 667 & 14.1 & 14 \\
0.98 & 3,730 & 994 & 21.0 & 20 \\
\textbf{0.95} & \textbf{2,777} & \textbf{1,947} & \textbf{41.2} & \textbf{73} \\
0.92 & 2,206 & 2,518 & 53.3 & 83 \\
0.90 & 1,946 & 2,778 & 58.8 & 83 \\
\bottomrule
\end{tabular}
\end{center}

At the chosen threshold of 0.95, \textbf{zero} clusters span more than one class,
indicating no cross-class over-merging.

\begin{figure}[htbp]
\centering
\includegraphics[width=\linewidth]{cluster_grids.png}
\caption{The six largest clusters. Each row is one photograph reproduced under different
colour transformations. The top row contains 73 copies of a single image. This is visual
confirmation that the clustering is identifying genuine duplicates rather than merely
similar photographs.}
\end{figure}

\subsection{What the augmentation actually was}

The grids make the mechanism plain: the author's augmentation was \textbf{colour
jitter}, not geometric transformation. This is precisely why a difference hash failed.
Greyscale luminance is computed as $L = 0.299R + 0.587G + 0.114B$; a hue shift changes
the channel mix, which changes $L$ unevenly across the frame, which flips the gradient
comparisons that constitute the hash --- while leaving the photograph's content
identical.

A difference hash is sensitive to exactly what the augmentation altered and blind to what
it preserved. It was the worst available tool for this particular dataset.

\subsection{Consequence for Submission 2}

\begin{center}
\begin{tabular}{@{}lr@{}}
\toprule
Test images, Submission 2 split & 721 \\
With a near-duplicate in training & \textbf{331} \\
\textbf{Contamination rate} & \textbf{45.9\%} \\
Genuinely held out & 390 (54\%) \\
\bottomrule
\end{tabular}
\end{center}

\subsection{An honest limit on this conclusion}

Contamination is proven and large. It does \textbf{not}, on its own, explain Submission
2's accuracy. With only two errors in 721 predictions, accuracy on the 390 genuinely
uncontaminated test images was between 0.9949 and 1.0000 whichever way the two errors
fell. Something beyond duplicate leakage was also inflating that result --- plausibly
session correlation (photographs from the same match sharing stadium, kit and lighting),
or a threshold that remains conservative, or a task that is simply easier than
Submission 1 assumed.

The present submission's drop to 0.8514 on the clean split is consistent with the
leakage having been substantial, but it cannot isolate the cause, because it also changed
architecture and training regime. The eleven-minute experiment described in Section 5
would isolate it.

\section{Conclusion}

\begin{enumerate}
\item A 1.2 M-parameter network trained from random initialisation reaches
  \textbf{0.8515 macro F1} on a leakage-controlled split of this dataset.
\item It is limited by its training budget rather than by capacity or overfitting, and
  would improve with a longer schedule.
\item The per-class difficulty ordering matches Submission 2's mean ordering, with the
  same easiest and hardest classes, and is best read as a property of the data.
\item The hardest confusion is \texttt{drive} $\leftrightarrow$
  \texttt{legglance-flick}, confirming Submission 1's visual prediction and correcting
  Submission 2's contaminated measurement.
\item Submission 4 has 14.8 pp of headroom against a 3.1 pp detection threshold. The
  comparison is viable, and would not have been without the audit.
\end{enumerate}

\end{document}
